% Section: K-Scaling — Hybrid vs. Disjoint LinUCB
% Include from the main paper via % Section: K-Scaling — Hybrid vs. Disjoint LinUCB
% Include from the main paper via % Section: K-Scaling — Hybrid vs. Disjoint LinUCB
% Include from the main paper via % Section: K-Scaling — Hybrid vs. Disjoint LinUCB
% Include from the main paper via \input{../experiments/05_figure/section_k_scaling_results}

\subsubsection{Ablation: Sample Efficiency of Family Parameter Sharing}
\label{sec:k_scaling}

We isolate the contribution of Hybrid LinUCB's family sharing by comparing it against Disjoint LinUCB (independent models) at $K{\in}\{5, 10, 20\}$. Both policies start tabula rasa, training on 888 prompts. We use data-driven family assignments via within-provider tetrachoric correlation ($r_{\text{tet}} \geq 0.6$) to correct for binary reward base-rate attenuation, avoiding naive syntactic groupings (e.g., separating uncorrelated \texttt{gpt-4-turbo} and \texttt{gpt-4.1}).

\begin{table}[h]
\centering
\caption{K-Scaling: Hybrid vs.\ Disjoint LinUCB. Holdout mean $\pm$ 95\% CI ($n{=}20$).}
\label{tab:k_scaling}
\small
\resizebox{\columnwidth}{!}{%
\begin{tabular}{@{}llcccr@{}}
\toprule
\textbf{K} & \textbf{Policy} & \textbf{Families} & \textbf{Holdout Reward} & \textbf{$p$-value} & \textbf{CI ratio} \\
\midrule
\multirow{2}{*}{5}
& Hybrid  & 4 & 0.966 $\pm$ 0.006 & \multirow{2}{*}{0.008**} & \multirow{2}{*}{3.2$\times$} \\
& Disjoint & --- & 0.937 $\pm$ 0.020 & & \\
\midrule
\multirow{2}{*}{10}
& Hybrid  & 6 & 0.958 $\pm$ 0.013 & \multirow{2}{*}{0.568} & \multirow{2}{*}{1.2$\times$} \\
& Disjoint & --- & 0.964 $\pm$ 0.016 & & \\
\midrule
\multirow{2}{*}{20}
& Hybrid  & 11 & 0.959 $\pm$ 0.007 & \multirow{2}{*}{0.320} & \multirow{2}{*}{3.1$\times$} \\
& Disjoint & --- & 0.946 $\pm$ 0.022 & & \\
\bottomrule
\end{tabular}%
}
\end{table}

\input{../experiments/05_figure/figure_5_caption}

\paragraph{Results \& Implications.}
Table~\ref{tab:k_scaling} and Figure~\ref{fig:k_scaling} show Hybrid LinUCB produces significantly more stable convergence. At $K{=}5$ and $K{=}20$, Hybrid's 95\% CI is $3\times$ narrower. This stability advantage arises because shared $\betavec_F$ parameters pool observations, providing a well-estimated baseline when per-model data is scarce ($44$~obs/model at $K{=}20$). 

For practitioners, this stability gain matters more than the mean gain: compressing the failure tail ensures predictable user experience. Furthermore, data-driven family assignment correctly groups models by empirical reward similarity rather than flawed syntactic naming conventions.


\subsubsection{Ablation: Sample Efficiency of Family Parameter Sharing}
\label{sec:k_scaling}

We isolate the contribution of Hybrid LinUCB's family sharing by comparing it against Disjoint LinUCB (independent models) at $K{\in}\{5, 10, 20\}$. Both policies start tabula rasa, training on 888 prompts. We use data-driven family assignments via within-provider tetrachoric correlation ($r_{\text{tet}} \geq 0.6$) to correct for binary reward base-rate attenuation, avoiding naive syntactic groupings (e.g., separating uncorrelated \texttt{gpt-4-turbo} and \texttt{gpt-4.1}).

\begin{table}[h]
\centering
\caption{K-Scaling: Hybrid vs.\ Disjoint LinUCB. Holdout mean $\pm$ 95\% CI ($n{=}20$).}
\label{tab:k_scaling}
\small
\resizebox{\columnwidth}{!}{%
\begin{tabular}{@{}llcccr@{}}
\toprule
\textbf{K} & \textbf{Policy} & \textbf{Families} & \textbf{Holdout Reward} & \textbf{$p$-value} & \textbf{CI ratio} \\
\midrule
\multirow{2}{*}{5}
& Hybrid  & 4 & 0.966 $\pm$ 0.006 & \multirow{2}{*}{0.008**} & \multirow{2}{*}{3.2$\times$} \\
& Disjoint & --- & 0.937 $\pm$ 0.020 & & \\
\midrule
\multirow{2}{*}{10}
& Hybrid  & 6 & 0.958 $\pm$ 0.013 & \multirow{2}{*}{0.568} & \multirow{2}{*}{1.2$\times$} \\
& Disjoint & --- & 0.964 $\pm$ 0.016 & & \\
\midrule
\multirow{2}{*}{20}
& Hybrid  & 11 & 0.959 $\pm$ 0.007 & \multirow{2}{*}{0.320} & \multirow{2}{*}{3.1$\times$} \\
& Disjoint & --- & 0.946 $\pm$ 0.022 & & \\
\bottomrule
\end{tabular}%
}
\end{table}

% Figure 5: K-Scaling Experiment — Hybrid vs. Disjoint LinUCB

\begin{figure*}[t]
\centering
\includegraphics[width=0.95\textwidth]{figures/figure5_k_scaling.png}
\caption{\textbf{Family parameter sharing reduces convergence variance in the production router.}
Greedy holdout reward (750 prompts) at 25-step intervals during training on 533 prompts via the full Corralling router (warmup + tabula rasa experts).
Columns show portfolios of $K \in \{5, 10\}$ models; bands are 95\% CIs ($n{=}20$ seeds).
Family assignments use within-provider tetrachoric correlation ($r_{\text{tet}} \geq 0.6$) computed on the training set only.
Both portfolios include multi-member providers so shared families form at every $K$: 2 shared families at $K{=}5$ (4/5 models) and 3 at $K{=}10$ (6/10 models).
Hybrid produces consistently narrower CIs---$1.8\times$ at $K{=}5$ ($\pm 0.018$ vs.\ $\pm 0.033$) and $2.1\times$ at $K{=}10$ ($\pm 0.017$ vs.\ $\pm 0.035$)---though mean differences are not significant at the $0.05$ level ($p{=}0.078$ and $p{=}0.160$, respectively).
Arm-pull allocation is highly non-uniform (e.g.\ median 47, range 0--120 at $K{=}10$), confirming that per-model data scarcity is driven by bandit exploitation, not uniform assignment.}
\Description{Two-panel line plot comparing Hybrid and Disjoint holdout reward over training steps at K=5 and K=10, with 95 percent confidence bands showing Hybrid producing narrower intervals at both portfolio sizes.}
\label{fig:k_scaling}
\end{figure*}


\paragraph{Results \& Implications.}
Table~\ref{tab:k_scaling} and Figure~\ref{fig:k_scaling} show Hybrid LinUCB produces significantly more stable convergence. At $K{=}5$ and $K{=}20$, Hybrid's 95\% CI is $3\times$ narrower. This stability advantage arises because shared $\betavec_F$ parameters pool observations, providing a well-estimated baseline when per-model data is scarce ($44$~obs/model at $K{=}20$). 

For practitioners, this stability gain matters more than the mean gain: compressing the failure tail ensures predictable user experience. Furthermore, data-driven family assignment correctly groups models by empirical reward similarity rather than flawed syntactic naming conventions.


\subsubsection{Ablation: Sample Efficiency of Family Parameter Sharing}
\label{sec:k_scaling}

We isolate the contribution of Hybrid LinUCB's family sharing by comparing it against Disjoint LinUCB (independent models) at $K{\in}\{5, 10, 20\}$. Both policies start tabula rasa, training on 888 prompts. We use data-driven family assignments via within-provider tetrachoric correlation ($r_{\text{tet}} \geq 0.6$) to correct for binary reward base-rate attenuation, avoiding naive syntactic groupings (e.g., separating uncorrelated \texttt{gpt-4-turbo} and \texttt{gpt-4.1}).

\begin{table}[h]
\centering
\caption{K-Scaling: Hybrid vs.\ Disjoint LinUCB. Holdout mean $\pm$ 95\% CI ($n{=}20$).}
\label{tab:k_scaling}
\small
\resizebox{\columnwidth}{!}{%
\begin{tabular}{@{}llcccr@{}}
\toprule
\textbf{K} & \textbf{Policy} & \textbf{Families} & \textbf{Holdout Reward} & \textbf{$p$-value} & \textbf{CI ratio} \\
\midrule
\multirow{2}{*}{5}
& Hybrid  & 4 & 0.966 $\pm$ 0.006 & \multirow{2}{*}{0.008**} & \multirow{2}{*}{3.2$\times$} \\
& Disjoint & --- & 0.937 $\pm$ 0.020 & & \\
\midrule
\multirow{2}{*}{10}
& Hybrid  & 6 & 0.958 $\pm$ 0.013 & \multirow{2}{*}{0.568} & \multirow{2}{*}{1.2$\times$} \\
& Disjoint & --- & 0.964 $\pm$ 0.016 & & \\
\midrule
\multirow{2}{*}{20}
& Hybrid  & 11 & 0.959 $\pm$ 0.007 & \multirow{2}{*}{0.320} & \multirow{2}{*}{3.1$\times$} \\
& Disjoint & --- & 0.946 $\pm$ 0.022 & & \\
\bottomrule
\end{tabular}%
}
\end{table}

% Figure 5: K-Scaling Experiment — Hybrid vs. Disjoint LinUCB

\begin{figure*}[t]
\centering
\includegraphics[width=0.95\textwidth]{figures/figure5_k_scaling.png}
\caption{\textbf{Family parameter sharing reduces convergence variance in the production router.}
Greedy holdout reward (750 prompts) at 25-step intervals during training on 533 prompts via the full Corralling router (warmup + tabula rasa experts).
Columns show portfolios of $K \in \{5, 10\}$ models; bands are 95\% CIs ($n{=}20$ seeds).
Family assignments use within-provider tetrachoric correlation ($r_{\text{tet}} \geq 0.6$) computed on the training set only.
Both portfolios include multi-member providers so shared families form at every $K$: 2 shared families at $K{=}5$ (4/5 models) and 3 at $K{=}10$ (6/10 models).
Hybrid produces consistently narrower CIs---$1.8\times$ at $K{=}5$ ($\pm 0.018$ vs.\ $\pm 0.033$) and $2.1\times$ at $K{=}10$ ($\pm 0.017$ vs.\ $\pm 0.035$)---though mean differences are not significant at the $0.05$ level ($p{=}0.078$ and $p{=}0.160$, respectively).
Arm-pull allocation is highly non-uniform (e.g.\ median 47, range 0--120 at $K{=}10$), confirming that per-model data scarcity is driven by bandit exploitation, not uniform assignment.}
\Description{Two-panel line plot comparing Hybrid and Disjoint holdout reward over training steps at K=5 and K=10, with 95 percent confidence bands showing Hybrid producing narrower intervals at both portfolio sizes.}
\label{fig:k_scaling}
\end{figure*}


\paragraph{Results \& Implications.}
Table~\ref{tab:k_scaling} and Figure~\ref{fig:k_scaling} show Hybrid LinUCB produces significantly more stable convergence. At $K{=}5$ and $K{=}20$, Hybrid's 95\% CI is $3\times$ narrower. This stability advantage arises because shared $\betavec_F$ parameters pool observations, providing a well-estimated baseline when per-model data is scarce ($44$~obs/model at $K{=}20$). 

For practitioners, this stability gain matters more than the mean gain: compressing the failure tail ensures predictable user experience. Furthermore, data-driven family assignment correctly groups models by empirical reward similarity rather than flawed syntactic naming conventions.


\subsubsection{Ablation: Sample Efficiency of Family Parameter Sharing}
\label{sec:k_scaling}

We isolate the contribution of Hybrid LinUCB's family sharing by comparing it against Disjoint LinUCB (independent models) at $K{\in}\{5, 10, 20\}$. Both policies start tabula rasa, training on 888 prompts. We use data-driven family assignments via within-provider tetrachoric correlation ($r_{\text{tet}} \geq 0.6$) to correct for binary reward base-rate attenuation, avoiding naive syntactic groupings (e.g., separating uncorrelated \texttt{gpt-4-turbo} and \texttt{gpt-4.1}).

\begin{table}[h]
\centering
\caption{K-Scaling: Hybrid vs.\ Disjoint LinUCB. Holdout mean $\pm$ 95\% CI ($n{=}20$).}
\label{tab:k_scaling}
\small
\resizebox{\columnwidth}{!}{%
\begin{tabular}{@{}llcccr@{}}
\toprule
\textbf{K} & \textbf{Policy} & \textbf{Families} & \textbf{Holdout Reward} & \textbf{$p$-value} & \textbf{CI ratio} \\
\midrule
\multirow{2}{*}{5}
& Hybrid  & 4 & 0.966 $\pm$ 0.006 & \multirow{2}{*}{0.008**} & \multirow{2}{*}{3.2$\times$} \\
& Disjoint & --- & 0.937 $\pm$ 0.020 & & \\
\midrule
\multirow{2}{*}{10}
& Hybrid  & 6 & 0.958 $\pm$ 0.013 & \multirow{2}{*}{0.568} & \multirow{2}{*}{1.2$\times$} \\
& Disjoint & --- & 0.964 $\pm$ 0.016 & & \\
\midrule
\multirow{2}{*}{20}
& Hybrid  & 11 & 0.959 $\pm$ 0.007 & \multirow{2}{*}{0.320} & \multirow{2}{*}{3.1$\times$} \\
& Disjoint & --- & 0.946 $\pm$ 0.022 & & \\
\bottomrule
\end{tabular}%
}
\end{table}

% Figure 5: K-Scaling Experiment — Hybrid vs. Disjoint LinUCB

\begin{figure*}[t]
\centering
\includegraphics[width=0.95\textwidth]{figures/figure5_k_scaling.png}
\caption{\textbf{Family parameter sharing reduces convergence variance in the production router.}
Greedy holdout reward (750 prompts) at 25-step intervals during training on 533 prompts via the full Corralling router (warmup + tabula rasa experts).
Columns show portfolios of $K \in \{5, 10\}$ models; bands are 95\% CIs ($n{=}20$ seeds).
Family assignments use within-provider tetrachoric correlation ($r_{\text{tet}} \geq 0.6$) computed on the training set only.
Both portfolios include multi-member providers so shared families form at every $K$: 2 shared families at $K{=}5$ (4/5 models) and 3 at $K{=}10$ (6/10 models).
Hybrid produces consistently narrower CIs---$1.8\times$ at $K{=}5$ ($\pm 0.018$ vs.\ $\pm 0.033$) and $2.1\times$ at $K{=}10$ ($\pm 0.017$ vs.\ $\pm 0.035$)---though mean differences are not significant at the $0.05$ level ($p{=}0.078$ and $p{=}0.160$, respectively).
Arm-pull allocation is highly non-uniform (e.g.\ median 47, range 0--120 at $K{=}10$), confirming that per-model data scarcity is driven by bandit exploitation, not uniform assignment.}
\Description{Two-panel line plot comparing Hybrid and Disjoint holdout reward over training steps at K=5 and K=10, with 95 percent confidence bands showing Hybrid producing narrower intervals at both portfolio sizes.}
\label{fig:k_scaling}
\end{figure*}


\paragraph{Results \& Implications.}
Table~\ref{tab:k_scaling} and Figure~\ref{fig:k_scaling} show Hybrid LinUCB produces significantly more stable convergence. At $K{=}5$ and $K{=}20$, Hybrid's 95\% CI is $3\times$ narrower. This stability advantage arises because shared $\betavec_F$ parameters pool observations, providing a well-estimated baseline when per-model data is scarce ($44$~obs/model at $K{=}20$). 

For practitioners, this stability gain matters more than the mean gain: compressing the failure tail ensures predictable user experience. Furthermore, data-driven family assignment correctly groups models by empirical reward similarity rather than flawed syntactic naming conventions.
